\section{Dual identities, normal phases and actual domains}\label{sec:duals}
For $g,h\in H$ let $\Theta(g\otimes h)$ be the rank-one operator
\[
 \Theta(g\otimes h)v:=g\tau(hv).
\]
The tracial pairing is defined for products of two $L^2$ vectors.
The map $\Theta$ is a unitary identification of $H\bar\otimes H$ with $\HS(H)$, the Hilbert--Schmidt operators for the usual trace $\Tr$ on $\mathcal B(H)$.
It takes algebraic tensors to $\F(H)$ and satisfies
\begin{equation}\label{eq:theta}
 \Theta(a\cdot q\cdot b)=a\Theta(q)b,\qquad
 \Theta(q)\Omega=(\id\otimes\tau)q.
\end{equation}
These conventions agree with the tracial realization used in \cite{Voi98,ACSY24}.

\begin{lemma}\label{lem:phase}
Suppose $\mathsf A(X)$ holds and each $X_j$ is a bounded normal operator.
Every marginal spectral law $\mu_j$ is diffuse, and the derivations are uniquely determined by
\begin{equation}\label{eq:phase}
 \partial_iX_j^*=\delta_{ij}k_j,\qquad
 k_j(z,w):=\frac{\bar z-\bar w}{z-w}\quad(z\ne w).
\end{equation}
Here $k_j$ belongs to the marginal subspace $L^2(\mu_j^2)$ of the coarse Hilbert space; its value on the diagonal may be set to zero.
\end{lemma}
\begin{proof}
By \cite[Proposition~3.5]{ACSY24}, $\Delta(X)=d$.
If $e=1_{\{z\}}(X_j)$ were nonzero, the rank-one projection onto $\C e\Omega$ would commute with $\rho(X_j)$, contradicting \eqref{eq:cycles}.
Differentiate $X_jX_j^*=X_j^*X_j$ to obtain
\[
 (\lambda(X_j)\otimes1-1\otimes\rho(X_j))\partial_iX_j^*
 =\delta_{ij}(X_j^*\otimes1-1\otimes X_j^*).
\]
The spectral diagonal of the operator on the left is zero on the entire coarse Hilbert space.
Indeed, the two marginal spectral actions form a normal representation of the product abelian algebra, whose faithful product trace gives the diagonal mass zero by diffuseness.
Consequently the operator is injective.
The vector in \eqref{eq:phase} solves this equation, and uniqueness follows.
\end{proof}

We shall also need a version of this conclusion starting from dual operators.
\begin{lemma}\label{lem:commutation}
Let $x\in\M$ be a bounded normal operator and let $E:\M\to H$ be bounded for the operator norm on $\M$.
If $[E,x]=P_\Omega$, then the law $\mu$ of $x$ is diffuse and
\[
 [E,x^*]=\Theta(k_x)|_{\M\Omega},\qquad
 k_x(z,w):=(\bar z-\bar w)/(z-w).
\]
If $[E,x]=0$, then $[E,x^*]=0$ without a diffuseness assumption.
\end{lemma}
\begin{proof}
For the homogeneous assertion the map
\[
 z\longmapsto e^{zx^*}Ee^{-zx^*}
 =e^{zx^*-\bar z x}Ee^{-zx^*+\bar z x}
\]
is entire with values in $\mathcal B(\M,H)$ and uniformly bounded.
The equality uses the assumed commutation with $x$, and the two displayed multipliers are unitary.
Scalar Liouville applied through continuous linear functionals gives the assertion by differentiation at zero.

In the inhomogeneous case a spectral atom $e$ at $\lambda$ would satisfy
$(x-\lambda)E(e)=-\tau(e)\Omega$.
Multiplying by $e$ gives a contradiction.
Put $K:=\Theta(k_x)$ and $F:=[E,x^*]-K|_\M$.
The kernel identity $[K,x]=[P_\Omega,x^*]$ and Jacobi's identity show that $[F,x]=0$.
The homogeneous assertion gives $[F,x^*]=0$.

Partition the spectrum into finitely many cells of diameter at most $\varepsilon$ with $\mu$-null boundaries, and let $e_r$ be their spectral projections.
Continuous cutoffs from inside a cell and its complement converge strongly to $e_r$ and $1-e_r$.
Since $F$ commutes with continuous functions of $x$, applying these cutoffs to the fixed vectors $F(e_ra)$ and $F(a)$ proves
\[
 F(e_ra)=e_rF(a)\qquad(a\in\M).
\]
No limit is passed through $F$.
Choose $\lambda_r$ in each cell and put $x_\varepsilon:=\sum_r\lambda_re_r$ and $E_\varepsilon:=\sum_re_rEe_r$ as maps $\M\to H$.
Averaging $u_\sigma E u_\sigma$ over the unitaries $u_\sigma:=\sum_r\sigma_re_r$, $\sigma_r\in\{1,-1\}$, gives $\|E_\varepsilon\|\le\|E\|$.
Since $E_\varepsilon$ commutes with $x_\varepsilon^*$,
\[
 \|F\|_{\mathcal B(\M,H)}
 \le2\varepsilon\|E\|+
 \left\|\sum_re_rKe_r\right\|_{\HS}.
\]
The last squared norm is at most $\mu^2\{|z-w|\le\varepsilon\}$ and tends to zero.
This proves the inhomogeneous assertion.
\end{proof}

\begin{proof}[Proof of $\mathsf S\Rightarrow\mathsf A$ in Theorem~\ref{thm:normal}]
Let $D_i$ witness $\mathsf S$ and put $\eta_i:=D_i\Omega$.
Each $D_i$ is closable, and its closure is defined on $\M\Omega$.
The closed graph theorem applied to $a\mapsto\overline{D_i}(a\Omega)$ from the Banach space $\M$ to $H$ makes $D_i|_{\M\Omega}$ bounded for the operator norm.
Set
\[
 E_i(a\Omega):=-D_i(a\Omega)+a\eta_i.
\]
Then $E_i\Omega=0$ and $[E_i,X_j]=\delta_{ij}P_\Omega$.
Lemma~\ref{lem:commutation} determines the adjoint-generator commutators as the normal phases.
By expanding commutators, $[E_i,p]$ is the restriction of a Hilbert--Schmidt operator for every $p\in\A_X$.
Define $\partial_i p:=\Theta^{-1}([E_i,p])$.
The extension from $\M\Omega$ is unique by density, so this definition respects every represented relation.
Equation~\eqref{eq:theta} proves the outer Leibniz rule and gives
\[
 (\id\otimes\tau)\partial_i p=E_i(p\Omega).
\]
Testing at the vacuum now yields
\[
 \langle\Omega\otimes\Omega,\partial_i p\rangle
 =\langle-D_i^*\Omega+\eta_i^*,p\Omega\rangle.
\]
This is the actual Hilbert adjoint condition and proves \eqref{eq:strong-vacuum}.
Uniqueness follows from the injective-gap argument of Lemma~\ref{lem:phase}.
\end{proof}

\begin{lemma}\label{lem:graph}
Let $\A\subset\M$ be a weakly dense unital star algebra, and let $\partial:\A\to H\bar\otimes H$ be an outer derivation with actual vacuum $\xi$.
Suppose its right and left slices $E$ and $L$ satisfy
\[
 \|Ep\|_2\le C_E\|p\|,\qquad \|Lp\|_2\le C_L\|p\|.
\]
The operator $E^\circ$ with domain $\A\Omega$ is closable, and
\begin{equation}\label{eq:slice-adjoint}
 (E^\circ)^*(p\Omega)=p\xi-(L(p^*))^*\qquad(p\in\A).
\end{equation}
Both $\overline{E^\circ}$ and $(E^\circ)^*$ are defined on $\M\Omega$.
Any bounded-generator commutator identity for $E^\circ$ passes to $\M\Omega$ for this closure.
\end{lemma}
\begin{proof}
The Leibniz rule gives
\[
 (\tau\otimes\tau)\partial(ab)=\tau(L(a)b)+\tau(aE(b)).
\]
Use the vacuum pairing on $p^*q$ to obtain \eqref{eq:slice-adjoint}.
Thus the adjoint has dense domain and
\[
 \|(E^\circ)^*(p\Omega)\|_2\le(\|\xi\|_2+C_L)\|p\|.
\]
For $a\in\M$, Kaplansky density followed by norm approximation supplies $p_\alpha\in\A$ with $p_\alpha\Omega\to a\Omega$ and $\sup_\alpha\|p_\alpha\|\le\|a\|+\varepsilon$.
The two image nets are bounded in $H$.
Pass to a common weakly convergent subnet and use weak closedness of both closed graphs.
It follows that $a\Omega$ belongs to both actual domains.
For a commutator, include the approximants $X_jp_\alpha$ in the same argument; the bounded commutator term identifies the graph limit.
\end{proof}

We record separately why the weakest full-domain condition still gives maximal $\Delta$.
\begin{proposition}\label{prop:bare-maximal}
For an arbitrary bounded tuple, $\mathsf F(X)$ implies $\Delta(X)=d$.
\end{proposition}
\begin{proof}
Use the equivalent left-multiplication form of \eqref{eq:cycles}, obtained by adjoints and tracial conjugation.
Suppose $\sum_i[T_i,X_i]=0$ with $T_i\in\F(H)$.
Choose a bounded regular $r\in\M$ such that every vector in the finite-dimensional space $\sum_i\Ran T_i$ becomes an element of $\M\Omega$ after right multiplication by $r$.
Such an $r$ exists by clearing a finite family of affiliated right denominators.
For $c,e\in\M$, put $Q:=\rho(r)$ and $S_i:=\rho(c)QT_i\rho(e)$.
All right multipliers commute with the left $X_i$, so the $S_i$ form another cycle and their ranges lie in $\M\Omega$.
Each $D_jS_i$ is an everywhere-defined bounded finite-rank operator: express its range through a finite basis of $\Ran S_i$ and apply the algebraic map $D_j$ to that basis.
Applying $D_j$ to the cycle, with the generator-first sign, gives
\[
 0=\sum_i[D_jS_i,X_i]+P_\Omega S_j.
\]
All operators in this equality are bounded and finite rank, so taking the usual trace gives
\[
 0=\Tr(P_\Omega S_j)
 =\langle\rho(c)^*\Omega,QT_j\rho(e)\Omega\rangle.
\]
The two test families are independently dense in $H$.
Hence $QT_j=0$, and injectivity of $Q$ gives $T_j=0$.
This proves maximality.
\end{proof}
